\documentclass[11pt]{article}
\usepackage[margin=1in]{geometry}
\usepackage{xcolor}
\usepackage{enumitem}
\usepackage{amsmath}
\usepackage{amssymb}
\usepackage{hyperref}
\hypersetup{colorlinks=true,linkcolor=black,urlcolor=blue}
\definecolor{revblue}{RGB}{0,0,180}
\definecolor{revred}{RGB}{180,0,0}
\newcommand{\loc}[1]{\textbf{[\textcolor{revblue}{Tracked PDF: #1}]}}
\setlist{nosep,leftmargin=1.4em}
\setlength{\parskip}{0.55em}
\setlength{\parindent}{0pt}

\begin{document}

\begin{center}
{\Large\textbf{Response to Reviewers}}\\[0.6em]
\textbf{Round 2 Revision}
\end{center}

\textbf{Manuscript:} ``Content-Conditioned Invariant Feature-Statistic
Augmentation: Unifying Robustness and Weight Averaging for Domain
Generalization'' (CIFA)\\
\textbf{Journal:} \textit{Electronics} (MDPI)\qquad
\textbf{Manuscript ID:} electronics-4654208\\
\textbf{Date:} 3 November 2026

\hrule
\vspace{0.4em}

Dear Editors and Reviewers,

We thank all three reviewers for the careful and constructive second-round
assessment. The remaining concerns are well focused: the readability of the
figures (Reviewer~1), the interpretation and positioning of the work together
with an editorial cleanup (Reviewer~2), and the mathematical justification,
terminology, and internal consistency of the manuscript (Reviewer~3). We have
addressed every point; no new training experiments were required.

All revisions are marked in the tracked manuscript
(\texttt{CIFA\_Electronics\_tracked.pdf}): additions in blue and deletions in
red with a strikethrough. The manuscript uses continuous line numbering
(right margin). Each response below is preceded by a location tag of the form
\loc{Page X, Lines Y--Z} that points to the corresponding revision in the
tracked PDF; a tag spanning two pages is written as
\loc{Page X, Line Y through Page X+1, Line Z}.

The revision comprises five substantive actions:
\begin{itemize}
\item \textbf{A1.} All seven figures redrawn as publication-ready vector
graphics---legends moved outside the plotting area, distinguishable
line styles/markers, enlarged fonts and insets, with no change to any data;
\item \textbf{A2.} Proposition~1 restructured so that the unconditional
distribution-level total-variation bounds are separated from the pointwise
posterior bound, the latter now stated under an explicit inter-model
likelihood-ratio condition ($\mathrm{TV}\le\tanh\alpha$);
\item \textbf{A3.} The wording made consistent across the paper:
non-significant differences are reported as statistically tied/on par, and
``equivalence'' is used only for the two-one-sided (TOST) test with the
equivalence margin stated; no superiority is claimed;
\item \textbf{A4.} Editorial cleanup, including removal of a duplicated
data-availability statement;
\item \textbf{A5.} A full consistency check of equations, tables,
cross-references, and terminology.
\end{itemize}

We reproduce each comment and respond point by point.

%==================================================================
\section*{Reviewer 1}

\textbf{Comment.} \textit{The remaining concern is the readability of the
figures. In several plots, portions of the curves are obscured by legends or
overlapping graphical elements, making it difficult to distinguish individual
traces and compare the results. The authors should carefully check all
figures, reposition legends outside the plotting areas where necessary, and use
clearly distinguishable line styles, colors, and markers. Enlarged panels or
zoomed-in insets would also help where curves overlap closely. Please ensure
that all curves, labels, and legends remain clearly readable at the final
publication size. I recommend acceptance after minor revision to address these
presentation issues.}

\loc{Figures 1--7, Pages 7, 14, 25, 29, 30, 33, 34}

\textbf{Response.} We agree, and we have redrawn \textbf{all seven figures}
(not only the overlapping panels) as publication-ready vector PDFs, while
leaving every data point unchanged. Each of the five requested actions has
been applied:

\begin{enumerate}[label=(\arabic*)]
\item \textbf{All figures checked and redrawn.} Figures~1--7 were rebuilt
from the identical source data with a uniform print specification;
\item \textbf{Legends moved outside the plotting area.} In-axis legends were
relocated to external margins (clearly visible in Figures~2, 6, and~7), so no
curve is occluded;
\item \textbf{Distinguishable styles and markers.} Marker size, line width,
and cap size were increased and styles differentiated for clear separation at
print size;
\item \textbf{Enlarged panels and zoomed insets.} The closely overlapping
coverage panel retains a zoomed inset (Figure~2b, zoom $3\times10$);
\item \textbf{Readability at publication size.} Each figure was re-laid-out at
its single- or double-column print dimensions and verified through a
print-size preview.
\end{enumerate}

The figure locations are: Figure~1 (Page~7), Figure~2 (Page~14), Figure~3
(Page~25), Figure~4 (Page~29), Figure~5 (Page~30), Figure~6 (Page~33), and
Figure~7 (Page~34). We confirm that the underlying values are identical to the
authoritative results and that no data were re-sampled.

%==================================================================
\section*{Reviewer 2}

\textbf{General opening.} We thank Reviewer~2 for the candid and careful
assessment, and for acknowledging that the presentation issues have been
addressed, that the claims have been appropriately revised, and that the null
and non-replicated results are reported transparently. The remaining concerns
concern the interpretation and positioning of the work rather than additional
experiments; we respond to each below.

\subsection*{(1) On the necessity of the distinctive components}
\loc{Page 33, Lines 1070--1074; Proposition 1, Pages 8--9, Lines 304--348}

We agree that the content--style decomposition is not a source of measured
accuracy: full-moment mixing is statistically indistinguishable from full
CIFA (paired $+0.23\pm0.75$, $p=0.53$), and removing weight averaging is the
dominant ablation ($2.42$ points). We would, however, distinguish accuracy
contribution from theoretical role. The decomposition is not intended to
improve accuracy; it is the structural condition that makes the
content-preservation bound admissible. By mixing only the style deviation and
fixing the content baseline, the construction guarantees
$\mathbb{E}[\mathbf{d}\mid Y]=0$ and thereby licenses the total-variation
bound of Proposition~1. Mixing the full moments does not satisfy this
condition and therefore cannot carry the same guarantee, even when its
accuracy is matched. The two constructions are tied in accuracy but not in the
guarantees they provide; empirical equivalence in accuracy does not establish
logical equivalence in the conditions each satisfies.

\subsection*{(2) On the value of the unification}
\loc{Pages 8--11, Lines 304--400}

We explicitly position the contribution as integration rather than a new
operator. The unification is not a concatenation of existing techniques: it
places feature-moment augmentation, convex-hull minimax, and weight averaging
under one causal-invariance principle and states the conditions under which
each applies. This yields testable predictions---content--style separability
as an enabling rather than sufficient condition, the behavior of the coverage
gap, and the second-order flat-minima approximation. Together with the
reproducible platform and the falsifiable boundary, this allows later methods
to be compared on a common footing and avoids repeated investment in
directions already shown to be ineffective. We do not claim a deployment
advantage; the practical significance lies in the unified understanding and
the well-characterized baseline.

\subsection*{(3) On the strength of the assumptions and the coverage term}
\loc{Page 8, Lines 304--322; Page 36, Lines 1171--1173}

The assumptions are now stated, tested, and falsifiable premises rather than
hidden conditions. Content--style separability is labeled an enabling premise,
and the trajectory--augmentation alignment is presented as an empirically
supported assumption rather than a theorem, with its limitations stated. The
target access required to evaluate the coverage-gap term is an inherent
property of target-free generalization, not a limitation unique to this work;
we demonstrate the behavior of both terms in a controlled moment-space setting
(Figure~2, Page~14) rather than presenting the term as computable without
target data. A strong assumption is not an incorrect one; the value of the
analysis is that its conditions are explicit and open to falsification.

\subsection*{(4) On weakening claims versus adding evidence}
\loc{Pages 21--23, Lines 795--828; Pages 8--11, Lines 304--400}

When the evidence does not support a stronger claim, weakening that claim is
the scientifically correct response; retaining an unsupported claim would be
the error. The revision also adds positive evidence---an independent
unified-code replication, a gate-pressure stress test, the alignment
verification, and a corrected two-one-sided equivalence analysis. This
evidence points to statistical equivalence and to boundary conditions rather
than to superiority. What was weakened were the claims that did not hold; the
claims that do hold---the unifying framework and the boundary
characterization---are retained, strengthened, and made more precise and more
readily falsifiable.

\subsection*{(5) Overall scientific significance}
\loc{Page 36, Lines 1168--1200}

The manuscript does not claim a performance breakthrough. Its scientific
significance is at the level of understanding: a unifying explanation, a
systematic validity boundary, and a falsifiable, reproducible benchmark; the
transparent reporting of null and non-replicated results is itself part of
that contribution. The only actionable editorial item has been completed (see
the editorial cleanup below). We respectfully ask that the work be judged on
what it actually establishes rather than on the overstated claims of the
original submission.

\subsection*{(6) Editorial cleanup (duplicated or partially edited text)}
\loc{Page 37, Lines 1209--1211}

We carried out a full editorial review. The one concrete duplication was the
data-availability statement, which appeared twice: a manually typeset
``\textbf{Data availability.}'' paragraph and the MDPI template
\texttt{\textbackslash dataavailability\{\}} command with identical content.
The manual paragraph has been removed and the template command retained, so
the statement now renders exactly once. The remaining items that can look like
duplication in a text extraction (redefined revision macros, deliberate
indentation and table alignment) were inspected individually and are not
duplicated content.

%==================================================================
\section*{Reviewer 3}

\textbf{General opening.} We sincerely thank Reviewer~3 for the detailed and
constructive assessment and for recognising the substantial improvement in the
mathematical formulation, experimental validation, statistical analysis, and
interpretation. We address the three limited revisions below.

\subsection*{R3-1. Mathematical justification for the conditional posterior bound}

\textbf{Comment.} \textit{Clarify the mathematical justification for
transferring the total-variation bound to the conditional content posteriors.
The data-processing inequality alone does not generally establish a pointwise
posterior bound under different conditioning observations. A more explicit
justification or an appropriately weakened statement would strengthen the
theoretical validity.}

\loc{Page 8, Line 304 through Page 9, Line 348; proofs on Page 10, Line 354
through Page 11, Line 400}

\textbf{Response.} The comment is correct, and we have restructured
Proposition~1 rather than defend the earlier statement. The data-processing
inequality controls the distributions (laws) of the posterior random
variables, not the posterior evaluated at each fixed observation; at rare
observations a small likelihood perturbation can be amplified. The proposition
is now stated at two clearly separated levels:

\begin{itemize}
\item \textbf{(L2-i) Exact case.} When $\tau=0$ and the partner source is
exogenous, the content posterior is preserved exactly;
\item \textbf{(L2-ii)(a) Distribution-level bounds (unconditional).} Without
further assumptions, the coupling gives the only unconditional statements: the
joint and observation-marginal laws
\[
\mathrm{TV}(P_{C,Z},P_{C,T(Z)})\le\tau,\qquad
\mathrm{TV}(p(Z),p(T(Z)))\le\tau;
\]
\item \textbf{(L2-ii)(b) Pointwise bound (conditional).} The pointwise
statement is retained only under an explicit inter-model likelihood-ratio
condition $|\log p(\mathbf{b},r,\mathbf{d}\mid c)-\log
p(\mathbf{b},r,\mathbf{d}_{aug}\mid c)|\le\alpha$, under which
$\mathrm{TV}(p(C\mid Z),p(C\mid T(Z)))\le\tanh\alpha$.
\end{itemize}

The KL corollary is now conditional on the same assumption,
$D_{\mathrm{KL}}\le\tanh(\alpha)/q_{\min}$, and the earlier
$\tau\log(M/p_{\min})$ and $\tau/q_{\min}$ pointwise bounds are removed. The
label-posterior drift $\delta$ is treated as a directly measured quantity
rather than added to $\tau$. The full derivations, including the
likelihood-ratio argument and the reverse-Pinsker inequality, are given in the
revised proof (Pages 10--11). Theorems~1 and~2 are unaffected because they use
the distribution-level total-variation quantity.

\subsection*{R3-2. Consistent distinction of equivalence, non-significance, and superiority}

\textbf{Comment.} \textit{The conclusions should consistently distinguish
between demonstrated statistical equivalence, non-significant differences, and
performance superiority. The conditional nature of the theoretical unification
should also remain clear.}

\loc{Page 1, Lines 17--19; Pages 21--23, Lines 795--828; Page 36,
Lines 1168--1200}

\textbf{Response.} We have applied one consistent terminology throughout the
twelve affected passages. Non-significant contrasts are reported as
\textbf{statistically tied} or \textbf{on par} (no significant difference);
the term \textbf{equivalence} is used only for the TOST procedure and always
with the equivalence margin stated; and no superiority is claimed. We also
corrected the TOST implementation to the standard two-one-sided procedure
using 90\% confidence intervals and recomputed every contrast. On PACS
($n=10$), at the one-point margin equivalence is confirmed only against RSC;
at the two-point margin it additionally holds against SWAD, CORAL, and Fish;
the CIFA--EFDM contrast is neither superior nor certified equivalent. These
wording changes appear at the abstract (Page~1), the introduction and
contributions (Pages~2--3), the experiments and captions (Pages~21--25), the
appendix (Pages 33--35), and the conclusions (Page~36). The conditional nature
of the unification is kept explicit in every theoretical passage.

\subsection*{R3-3. Clean manuscript and consistency of equations, tables, references, and terminology}

\textbf{Comment.} \textit{Provide a clean manuscript incorporating the
revisions and verify the consistency of equations, tables, cross-references,
and terminology.}

\loc{Whole manuscript; hyphenation at Page 36, Line 1168}

\textbf{Response.} A clean, fully accepted version
(\texttt{CIFA\_Electronics\_clean.pdf}, 39 pages) is provided alongside the
tracked version. We have verified that: every cross-reference resolves with no
dangling or duplicate labels; all 27 bibliography entries are cited; equations,
tables, and figures use automatic numbering with no hard-coded numbers; and the
terminology and notation are consistent. The clean manuscript renders entirely
in black with no revision marks. As part of this pass we also standardised the
hyphenation (e.g., ``content--style''). Both the clean and tracked manuscripts
compile with zero errors and zero overfull boxes, and their accepted content is
identical.

\vspace{0.6em}
We hope these revisions fully address the reviewers' concerns and we thank the
editorial team and reviewers once again for their time and guidance.

\vspace{0.8em}
Sincerely,\\
Qingsong Cai (on behalf of all authors)\\
\textit{School of Computer and Artificial Intelligence, Beijing Technology and
Business University}\\
\textit{Email: caiqs@btbu.edu.cn}

\end{document}
